

%% ===================================================================
\section{\vrev{C2 가우시안 불확실성 결과}}
\label{sec:result_c2}
%% ===================================================================

\textbf{C2} 조건은 C1의 관측성에 가우시안 불확실성($U_{\text{raw}}^{(G)} \sim N(3,1)$)을
추가한 것으로, 경꼬리(light-tail, 가우시안) 불확실성이 $\DRTO$에 미치는 효과를
관측하는 조건이다.
개별 시그모이드 바운딩 모델에서 $E_{O,\text{bnd}}$와 $E_{U,\text{bnd}}$가
각각 독립적으로 바운딩되므로,
$\DRTO = R_0 + E_{O,\text{bnd}} + E_{U,\text{bnd}}$의 구조가
관측성의 하향 효과와 불확실성의 상향 효과를 명확히 분리한다.

\figref{fig:gaussian_concept}\refpo{fig:gaussian_concept}{은}{는} 가우시안 불확실성의 분포 특성과
시그모이드 바운딩을 통한 $\DRTO$ 상향 조정 메커니즘을 개념적으로 도식화한 것이다.
\chg{패널 (a)는 평균 $\mu = 3.0$의 대칭적 확률밀도함수 $N(3,1)$을, 패널 (b)는 불확실성 수준에 따른
$\DRTO$의 상향 조정을 보이며, 시그모이드 바운딩에 의해 $R_0 < \DRTO < R_{\max}$가 보장되어 높은
$U$ 영역에서 자연스럽게 포화한다.}

\begin{figure}[htbp]
  \centering
  \begin{tikzpicture}
    % ── Panel (a): Gaussian N(3,1) PDF ──
    \begin{axis}[
      name=panelA,
      width=0.47\textwidth,
      height=6.5cm,
      xlabel={$U_{\text{raw}}$ (불확실성 수준, 시간)},
      ylabel={확률밀도},
      title={\textbf{(a)} 가우시안 불확실성 $U \sim N(3, 1)$},
      title style={font=\small\bfseries},
      xmin=-1, xmax=7,
      ymin=0, ymax=0.46,
      ymajorgrids=true,
      grid style={dashed, black},
      legend style={at={(0.02,0.98)}, anchor=north west, font=\scriptsize},
      ]
      % Gaussian PDF (approximated with smooth curve)
      \addplot[blue!70!black, thick, smooth, domain=-1:7, samples=100]
        {1/(sqrt(2*pi)*1)*exp(-0.5*((x-3)/1)^2)};
      \addlegendentry{$N(3, 1)$}

      % Mean line
      \draw[red!70!black, dashed, thick] (axis cs:3, 0) -- (axis cs:3, 0.40);
      \node[font=\scriptsize, red!70!black, fill=white, inner sep=1.5pt,
            rounded corners=1pt] at (axis cs:3, 0.44) {$\mu = 3.0$};

      % ±1σ lines
      \draw[black, dotted, thin] (axis cs:2, 0) -- (axis cs:2, 0.245);
      \draw[black, dotted, thin] (axis cs:4, 0) -- (axis cs:4, 0.245);
      \node[font=\tiny, black] at (axis cs:5.5, 0.26) {$\pm 1\sigma$: 68.3\%};
      \draw[-{Stealth}, black, thin] (axis cs:5.2, 0.25) -- (axis cs:4.1, 0.24);

      % Symmetry annotation
      \node[font=\tiny, black, align=center] at (axis cs:-0.2, 0.30) {대칭 분포\\(Skewness $= 0$)};
    \end{axis}

    % ── Panel (b): DRTO shift via sigmoid bounding ──
    \begin{axis}[
      at={(panelA.east)},
      anchor=west,
      xshift=18mm,
      width=0.47\textwidth,
      height=6.5cm,
      xlabel={$U_{\text{raw}}$ (불확실성 수준, 시간)},
      ylabel={$\DRTO$ (시간)},
      title={\textbf{(b)} 시그모이드 바운딩에 의한 불확실성 프리미엄},
      title style={font=\small\bfseries},
      xmin=0, xmax=8,
      ymin=5.0, ymax=25,
      ymajorgrids=true,
      grid style={dashed, black},
      legend style={at={(0.98,0.98)}, anchor=north east, font=\scriptsize},
      ]
      % DRTO = R0 + E_U,bnd = R0 + (Rmax-R0)·S(κ·U/6)
      % S(z) = 2/(1+exp(-z)) - 1
      \addplot[red!70!black, thick, smooth, domain=0:8, samples=100]
        {6.0 + 18.0*(2/(1+exp(-1.2*x/6)) - 1)};
      \addlegendentry{$R_0 + E_{U,\text{bnd}}$}

      % R0 baseline
      \draw[black, dashed, thick] (axis cs:0, 6.0) -- (axis cs:8, 6.0);
      \node[font=\tiny, black] at (axis cs:1.0, 6.5) {$R_0 = 6.0$h};

      % Rmax upper bound
      \draw[black, dotted, thin] (axis cs:0, 24.0) -- (axis cs:8, 24.0);
      \node[font=\tiny, black] at (axis cs:1.0, 24.5) {$R_{\max} = 24.0$h};

      % Mean point
      \addplot[only marks, mark=*, mark size=2.5pt, red!70!black]
        coordinates {(3.0, {6.0 + 18.0*(2/(1+exp(-1.2*3.0/6)) - 1)})};
      \node[font=\tiny, red!70!black, fill=white, inner sep=1pt, rounded corners=1pt,
            anchor=north west] at (axis cs:3.3, 9.0)
        {$U_{\text{raw}} = \mu\;(3.0)$};

      % Saturation annotation
      \node[font=\tiny, blue!70!black, align=center, fill=white, inner sep=1pt,
            rounded corners=1pt] at (axis cs:6.5, 19)
        {시그모이드 포화\\($< R_{\max}$)};
      \draw[-{Stealth}, blue!70!black, thin] (axis cs:6.5, 18.2) -- (axis cs:6.5, 16.5);
    \end{axis}
  \end{tikzpicture}
  \caption{가우시안 불확실성 개념도}
  \label{fig:gaussian_concept}
\end{figure}


%% -------------------------------------------------------------------
\subsection{\vrev{기술통계 및 효과 크기}}
\label{subsec:c2_descriptive}
%% -------------------------------------------------------------------

\chg{기술통계 측면에서,} C2 조건의 효율성 비율 평균은 $\bar{\eta}_{\text{C2}} = 1.682$,
$\mathrm{SD} = 0.615$로 나타났다.
$\DRTO$ 평균은 $10.094$시간으로서, 기준선 $R_0 = 6.0$시간 대비
$68.2\%$가 증가하였다.
C1의 $\bar{\eta} = 0.853$ 대비 $\erf{0.829}{0.830}$가 상승한 것으로,
불확실성의 도입이 $\eta$를 $1.0$ 이상으로
끌어올리는 \textbf{불확실성 프리미엄(uncertainty premium)} 효과를
확인하였다. 즉, 불확실성이 존재하는 환경에서 $\DRTO$는 기본 복구 시간
$R_0$보다 길어지며, 이는 복구 과정의 불확실성을 사전에 반영한
안전 마진(safety margin)으로 해석된다.

\chg{불확실성 프리미엄 비율을 보면,} 전체 $n = 10{,}000$개 표본 중 $\eta_{\text{C2}} > 1.0$인 비율,
즉 $\DRTO$가 기준선 $R_0$를 초과하는 비율은 $84.9\%$이다.
이는 가우시안 불확실성이 존재하는 환경에서 약 $85\%$의 시나리오가
기준선보다 긴 복구 시간을 필요로 함을 의미하며,
불확실성 프리미엄이 예외적 현상이 아닌 \textit{일반적 현상}임을 보여준다.
나머지 $15.1\%$는 관측성이 매우 높은($k_O \approx 1$, $O_{\text{raw}} \approx 1$) 조합에서
발생하는 것으로, 높은 관측성이 불확실성 프리미엄을 상쇄할 수 있는
조건을 나타낸다.

\chg{효과 크기 측면에서,} C1 대비 C2의 Cohen's $d = +1.871$으로서, \textbf{대효과(large effect)}에
해당한다(대효과 기준값 $|d|=0.80$의 약 $2.34$배).
방향이 양(+)인 것은 불확실성 추가가 $\eta$를 증가(복구 시간
연장)시킨다는 것을 나타내며, C0$\to$C1의 감소 방향($d = -1.180$)과
대조를 이룬다. 이러한 방향 반전은 관측성의 단축 효과와 불확실성의
연장 효과가 구조적으로 상반됨을 보여주며,
개별 시그모이드 바운딩 모델의 이론적 설계가 실험적으로 타당함을 입증한다.


%% -------------------------------------------------------------------
\subsection{\vrev{순서 준수 및 가설 검증}}
\label{subsec:c2_order}
%% -------------------------------------------------------------------

$n = 10{,}000$개 표본 중 $\eta_{\text{C1}} \leq \eta_{\text{C2}}$인
비율이 $99.87\%$로(\jrev{H3}-b), C2의 불확실성 프리미엄은 거의 모든
표본에서 관찰되었다. $0.13\%$의 예외는 관측성이 매우 높고($O_{\text{raw}} \approx 1$,
$k_O \approx 1$) 불확실성이 매우 낮은($U_{\text{raw}}^{(G)} \approx 0$) 극단적 조합에서
발생하는 것으로, 모델의 구조적 특성과 일치한다.
\jrev{H3}-b의 기준($\geq 99\%$)을 $99.87\%$로 충족한다.

\chg{C2의 순서 준수에 대한 H3-b의 검증 결론은 다음과 같다.} \mrev{H3-b는 ``$\eta_{\text{C1}} \leq \eta_{\text{C2}}$ 비율 $\geq 99\%$''를
기준으로 한다. 99.87\%가 충족되어 ($0.13\%$ 예외는 $O_{\text{raw}} \approx 1$
극단 영역), \textbf{H3-b는 채택된다}.}
